\section{Loop spaces, actions and sublevels}\label{sec:setup}

We fix notation for the geometric filtration and the three symmetry groups.
Except for the use of the classical two-sphere theorem, assume $n\ge3$ and
fix a smooth Riemannian metric $g$ on $S^n$. Put $d=n-1$ and
\[
 \Lambda=H^1(\mathbb R/\mathbb Z,S^n),\qquad
 M=\{\text{constant loops}\}\cong S^n.
\]
We use the strong $H^1$ topology. The energy, its square root, and length are
\[
 E(\gamma)=\int_0^1|\dot\gamma(t)|_g^2\dd t,
 \qquad F(\gamma)=\sqrt{E(\gamma)},
 \qquad \ell(\gamma)=\int_0^1|\dot\gamma(t)|_g\dd t.
\]
Then $\ell\le F$, with equality for constant-speed loops. Write
\[
 X_a=\{F<a\},\qquad \mathcal A_a=\{\ell<a\}\qquad(a\in\mathbb R).
\]
Both spaces are empty for $a\le0$; pairs relative to $M$ are used only
for $a>0$.
A positive cut $a$ is \emph{regular} if $a^2$ is a regular value of $E$.
At regular cuts we freely use the corresponding closed finite-dimensional
sublevel models, through the homotopy equivalences constructed in
Section~\ref{sec:local}. All maps between sublevels are the maps induced by
inclusion.

The orthogonal group $G=O(2)$ acts by changing the loop parameter. We single
out reflection $r$ and half-translation $\tau$:
\[
 (r\gamma)(t)=\gamma(-t),\qquad
 (\tau\gamma)(t)=\gamma(t+\tfrac12),\qquad
 R=\langle r\rangle,\quad K=\langle r,\tau\rangle.
\]
Thus $K\cong C_2\times C_2$. The fixed locus of $\tau$ consists of double
loops. The map
\[
 \delta:\Lambda\longrightarrow\Lambda^\tau,
 \qquad (\delta\gamma)(t)=\gamma(2t),
\]
is an $R$-equivariant homeomorphism, and $F(\delta\gamma)=2F(\gamma)$.
In particular, $(X_{2a})^\tau$ is the double of $X_a$.

All homology and cohomology have coefficients $\kk$. For an $H$-invariant
pair $(Y,Z)$, let
\[
 H_H^j(Y,Z)=H^j(EH\times_HY,EH\times_HZ;\kk).
\]
Ordinary groups have no group subscript. Negative-degree groups are zero,
and a total graded group always means the direct sum of its homogeneous
components. We use
\[
 B_a^j=H_R^j(X_a,M),\qquad
 B_\infty^j=H_R^j(\Lambda,M).
\]
The determinant character of $G$ defines $w\in H^1(BG;\kk)$, with the same
notation for its restrictions. Let $U\in H^2(BG;\kk)$ be the mod-two Euler
class of the standard plane. On $K$, let $v$ be the degree-one character
which is nonzero on $\tau$ and zero on $r$. Proposition~\ref{prop:gysin}
establishes the ring formulas and the natural sequences used below.

Evaluation at zero is $R$-equivariant for the trivial action on $S^n$.
If $\sigma\in H^n(S^n;\kk)$ is its nonzero class, write
\[
 e=\ev_0^*\sigma\in H_R^n(\Lambda;\kk).
\]
Cup product by $e$ acts on every relative reflection group. This operation
will distinguish a global pair of classes whose two members cannot first
appear at the same isolated critical crossing.
